%
% 14-analytisch.tex -- Potenzreihenentwicklung einer holomorphen Funktion
%
% (c) 2026 Prof Dr Andreas Müller
%
\subsection{Potenzreihenentwicklung einer holomorphen Funktion}
Falls für eine komplexe Funktion eine konvergente Potenzreihenentwicklung
existiert, dann ist die Funktion sogar unendlich oft stetig differnzierbar.
Ob auch die Umkehrung gilt, wissen wir noch nicht.
Die Existenz einer konvergenten Potenzreihe und die Existenz einer
komplexen Ableitung sind also nicht dasselbe.
Wir geben daher der Eigenschaft, eine konvergente Potenzreihe zu haben,
einen Namen.

\begin{definition}[analytisch]
Eine komplexe Funktion heisst im Punkt $z_0$ \emph{analytisch}, wenn sie eine
konvergente Potenzreihen im Punkt $z_0$ besitzt.
\index{analytisch}%
Eine reelle Funktion heisst im Punkt $x_0$ \emph{reell analytisch}, wenn sie
eine konvergente, reelle Potenzreihenentwicklung im Punkt $x_0$ besitzt.
\index{reell analytisch}%
\end{definition}

Falls $\gamma$ ein geschlossener Weg ist, der sich genau einmal
im Gegenuhrzeigersinn um den Punkt $a$ windet, lassen sich die
Integrale auf der rechten Seite von \eqref{buch:analytisch:eqn:reihe}
mit dem Satz~\ref{buch:integration:satz:ableitungn}
berechnet.
Wegen
\[
f^{(k)}(z)
=
\frac{k!}{2\pi i}
\oint_\gamma
\frac{f(x)}{(x-z)^{k+1}}
\,dx
\]
ergeben sie
\begin{align*}
f(z)
&=
\sum_{k=0}^{\infty}
\biggl(
\frac{1}{2\pi i}
\oint_\gamma
\frac{f(x)}
{(x-a)^{k+1}}
\,dx
\biggr)
(z-a)^k
=
\sum_{k=0}^\infty
\frac{f^{(k)}(a)}{k!}
(z-a)^k.
\end{align*}
Damit ist die gefundene Potenzreihe der Funktion $f$ die
Taylor-Reihe der Funktion $f$ im Punkt $a$.

\begin{satz}[holomorphe Funktionen sind analytisch]
\label{buch:analytisch:satz:holomorph-analytisch}
Eine holomorphe Funktion $f(z)$ ist analytisch, die Taylor-Reihe
von $f$ im Punkt $a$ konvergiert.
Der Konvergenzradius ist das Supremum der Radien von Kreisen um $a$,
die vollständig im Definitionsbereich von $f$ enthalten sind.
\end{satz}

Wie in der Einleitung zu diesem Abschnitt bemerkt wissen wir bereits,
dass analytische komplexe Funktionen auch holomorph sind, da eine
Potenzreihe gliedweise differenziert werden kann.
Der Satz zeigt jetzt, dass eine holomorphe Funktion immer auch
analytisch ist.
Dies ist eine Eigenschaft, die nicht für reelle analytische
Funktionen gilt, wie im nachfolgenden Abschnitt gezeigt wird.
